\chapter{Conclusões}\label{cap:Conclusões}

%\chapter{Discussão}\label{cap:discussão}

A presente pesquisa teve como objetivo analisar a seca extrema no Sudeste do Brasil durante o verão de 2013/2014, utilizando a decomposição dos modos normais tridimensionais para compreender a dinâmica atmosférica subjacente. Através da análise detalhada dos modos de ondas atmosféricas e da energia associada, foi possível identificar os principais mecanismos que impulsionaram a persistência das condições secas na região. A análise de precipitação e modos normais permitiu reconhecer o evento em três etapas. Etapa de iniciação, que abrange de 01 a 25 de dezembro de 2013 (ET1), Etapa intensa ou crítica, de 26 de dezembro de 2013 a 05 de fevereiro de 2014 (ET2), e etapa de decadência que vai de 06 a 28 de fevereiro de 2014 (ET3)   

A convecção profunda localizada no Pacífico Ocidental, desempenhou um papel crucial. Modulou a célula de Walker e gerou uma fonte eficiente de ondas de Rossby. Estas ondas, por sua vez, propagaram-se de forma guiada pelo jato subtropical em direção ao Atlântico Sul, onde reforçaram o padrão de bloqueio. A análise do vento zonal evidenciou um cisalhamento vertical intensificado na região do jato subtropical, favorável durante a etapa crítica (ET2), condição essencial para a conversão de energia de perturbações baroclínicas em modos barotrópicos de grande escala, sustentando assim a teleconexão tropical-extratropical.

A decomposição modal tridimensional provou ser uma ferramenta poderosa para isolar e quantificar a contribuição de diferentes tipos de ondas na dinâmica do evento. Ondas de Rossby (RH), foram as principais responsáveis pela configuração do padrão de bloqueio. A transição para uma fase positiva sobre o SEB durante a ET2 esteve diretamente associada à subsidência e à supressão da convecção. A energia associada a estas ondas concentrou-se nos modos verticais barotrópicos (m=1-3) e em números de onda zonal baixos (escalas planetárias), corroborando sua natureza de grande escala. 

A fase positiva das ondas de Kelvin (KW) predominante sobre a América do Sul, particularmente nos modos baroclínicos (m=7-9), contribuiu para uma dinâmica desfavorável à precipitação no SEB, possivelmente ligado a anomalias de circulação em grandes escalas.

Ondas de Gravidade (EIG/WIG) e Mistas (MRG) exibiram uma variabilidade espacial e temporal mais rápida, atuando como moduladoras da convecção em regiões como o Pacífico Ocidental. A sua energia, distribuída em modos baroclínicos, sugere uma ligação com fontes de calor convectivo e uma participação na transferência de energia vertical.

A análise energética revelou uma alteração significativa na partição de energia entre os modos verticais durante o evento. Verificou-se um claro predomínio da energia nos modos barotrópicos para as ondas de Rossby, enquanto as ondas de gravidade e de Kelvin apresentaram maior energia em modos baroclínicos. Adicionalmente, identificaram-se indícios de uma contribuição da estratosfera para a energia cinética e potencial disponível, especialmente nos modos baroclínicos superiores (e.g., m=8). Este achado sugere que processos de acoplamento vertical podem ter tido um papel modulador na persistência do evento, uma vez que a energia estratosférica mostrou-se relevante em regiões tropicais e subtropicais.

A abordagem por modos normais, utilizada neste estudo, não apenas corroborou os mecanismos já discutidos na literatura, como também contribuiu para o avanço na compreensão dos processos subjacentes ao evento. A decomposição modal permitiu quantificar a contribuição energética de diferentes tipos de ondas e apontou para uma possível influência de processos estratosféricos, ampliando o entendimento sobre o comportamento atmosférico durante a seca.

Os resultados sobre a energia regional mostram que os modos verticais de ordem superior, em conjunto, os modos 7, 8 e 9, são os principais responsáveis pela distribuição da energia durante as etapas de seca no sudeste do Brasil. Em particular, a energia potencial na ET2 predomina nas três regiões, indicando maior quantidade de energia disponível na troposfera durante este período. Este achado destaca a importância desses modos verticais na dinâmica atmosférica

Esses achados ressaltam a importância das interações entre as ondas atmosféricas e a contribuição estratosférica para a previsão de eventos climáticos extremos, como secas. Combinadas, essas abordagens fornecem uma visão detalhada dos mecanismos de grande escala envolvidos e podem aprimorar os modelos de previsão climática em futuras simulações. Além disso, os resultados têm implicações significativas para a modelagem de fenômenos climáticos extremos, como secas e inundações, e podem contribuir para o desenvolvimento de estratégias de mitigação e adaptação às mudanças climáticas, especialmente na América do Sul.
